
\chapter{Conclusão}
\label{chap:conclusao}

% Retomar brevemente o problema que motivou o desenvolvimento
% do DoseCerta e o objetivo geral do trabalho.

\section{Síntese do Trabalho}
\label{sec:sintese-trabalho}

% Resumir a proposta e o caminho seguido para desenvolver
% o sistema, destacando os principais subsistemas.

\section{Atendimento aos Objetivos}
\label{sec:conclusao-objetivos}

% Indicar em que medida o objetivo geral e os objetivos específicos
% foram alcançados, com base nos resultados e na discussão.
% Evitar afirmar que um objetivo foi atingido sem evidências.

\section{Principais Contribuições}
\label{sec:contribuicoes}

% Destacar as contribuições concretas do trabalho:
% desenvolvimento do protótipo, integração de componentes,
% implementação de funcionalidades ou aprendizados técnicos.
% Selecionar apenas contribuições efetivamente alcançadas.

\section{Limitações}
\label{sec:conclusao-limitacoes}

% Apresentar de forma concisa as limitações mais relevantes
% que permaneceram ao final do projeto.
% Não repetir toda a análise do Capítulo 16.

\section{Trabalhos Futuros}
\label{sec:trabalhos-futuros}

% Indicar melhorias e etapas futuras coerentes com as limitações:
% aperfeiçoamentos do protótipo, testes adicionais,
% avaliação com usuários ou desenvolvimento de funções pendentes.
% Priorizar propostas viáveis e justificadas.

% Encerrar com uma síntese do significado do trabalho
% e de sua contribuição para o problema investigado.
